\documentclass[12pt]{article}
\usepackage[utf8]{inputenc}
\usepackage[spanish]{babel}
\usepackage{geometry}
\usepackage{listings}
\usepackage{xcolor}
\usepackage{hyperref}

\geometry{a4paper, margin=1in}

% --- Configuración de estilo para el código Python ---
\definecolor{codegreen}{rgb}{0,0.6,0}
\definecolor{codegray}{rgb}{0.5,0.5,0.5}
\definecolor{codepurple}{rgb}{0.58,0,0.82}
\definecolor{backcolour}{rgb}{0.95,0.95,0.92}

\lstdefinestyle{mystyle}{
    backgroundcolor=\color{backcolour},   
    commentstyle=\color{codegreen},
    keywordstyle=\color{magenta},
    numberstyle=\tiny\color{codegray},
    stringstyle=\color{codepurple},
    basicstyle=\footnotesize\ttfamily,
    breakatwhitespace=false,         
    breaklines=true,                 
    captionpos=b,                    
    keepspaces=true,                 
    numbers=left,                    
    numbersep=5pt,                  
    showspaces=false,                
    showstringspaces=false,
    showtabs=false,                  
    tabsize=2
}
\lstset{style=mystyle}

\title{Tarea 1: Estructura del Proyecto y Módulo DNS}
\author{Prof. CESAR RODRIGUEZ}
\date{\today}

\begin{document}
\maketitle

\section{Introducción al Proyecto}

\textbf{Proyecto Semestral del Curso de Seguridad Informática}. A lo largo de las próximas 10 semanas, construiremos una herramienta de auditoría de seguridad en Python. El proyecto está estructurado en dos archivos principales que nos permitirán trabajar de forma modular y organizada.

\begin{itemize}
    \item \textbf{\texttt{auditoria.py}}: Será el punto de entrada de nuestro programa. Se encargará de gestionar los argumentos de la línea de comandos y de llamar a las funciones específicas que realizarán el trabajo.
    \item \textbf{\texttt{modulos.py}}: Será nuestra caja de herramientas. Contendrá todas las funciones que implementan las diferentes técnicas de auditoría (reconocimiento, escaneo, etc.).
\end{itemize}

A continuación se presenta el código base para ambos archivos.

\section{Archivo Principal: \texttt{auditoria.py}}
\subsection*{Explicación}
Este script es el orquestador de nuestra herramienta. Utiliza la librería \texttt{argparse} para crear una interfaz de línea de comandos profesional. Su función principal es interpretar qué módulos de auditoría desea ejecutar el usuario (por ejemplo, \texttt{--dns}, \texttt{--scan}) y llamar a las funciones correspondientes definidas en el archivo \texttt{modulos.py}, pasándoles el objetivo a auditar.

\subsection*{Código Fuente}
\begin{lstlisting}[language=Python, caption={auditoria.py}]
# auditoria.py
import argparse
import modulos  # Importamos nuestro archivo de modulos

def main():
    # --- Configuracion de Argumentos ---
    parser = argparse.ArgumentParser(
        description="Herramienta de Auditoria de Seguridad Informatica.",
        epilog="Ejemplo de uso: python auditoria.py example.com --dns --scan '1-1024'"
    )
    parser.add_argument("target", help="El objetivo (dominio o IP) a auditar.")
    
    # Argumentos para seleccionar los modulos a ejecutar
    parser.add_argument(
        "--dns", 
        action="store_true", 
        help="Ejecuta el modulo de reconocimiento DNS y WHOIS."
    )
    parser.add_argument(
        "--scan", 
        metavar="PUERTOS", 
        help="Ejecuta el escaner de puertos. E.g., '21,22,80' o '1-1024'."
    )

    args = parser.parse_args()
    target = args.target

    print(f"[*] Iniciando auditoria para el objetivo: {target}\n")

    # --- Logica para llamar a los modulos ---
    if args.dns:
        print("\n" + "="*20 + " MODULO: RECONOCIMIENTO DNS Y WHOIS " + "="*20)
        modulos.get_dns_info(target)
        print("="*65)

    if args.scan:
        print("\n" + "="*20 + " MODULO: ESCANER DE PUERTOS " + "="*25)
        open_ports = modulos.scan_ports(target, args.scan)
        print(f"\n[+] Escaneo de puertos finalizado. Se encontraron {len(open_ports)} puertos abiertos.")
        print("="*65)

    print("\n[*] Auditoria finalizada.")

if __name__ == "__main__":
    main()
\end{lstlisting}

\section{Archivo de Módulos: \texttt{modulos.py}}
\subsection*{Explicación}
Este archivo es el corazón funcional de nuestro proyecto. Aquí residirán todas las funciones que realizan las tareas de auditoría. Inicialmente, contiene una función genérica para el reconocimiento DNS y otra para el escaneo de puertos. A lo largo del semestre, cada estudiante añadirá sus propias funciones especializadas a este archivo, expandiendo las capacidades de nuestra herramienta.

\subsection*{Código Fuente}
\begin{lstlisting}[language=Python, caption={modulos.py}]
# modulos.py
import socket
import whois
import dns.resolver
from concurrent.futures import ThreadPoolExecutor

# ==============================================================================
# MODULO 1: RECONOCIMIENTO DNS Y WHOIS (BASE)
# ==============================================================================
def get_dns_info(domain):
    """
    Realiza consultas WHOIS y DNS para un dominio dado y muestra la informacion.
    """
    print(f"\n[+] Obteniendo informacion WHOIS para: {domain}")
    try:
        w = whois.whois(domain)
        print("  [-] Registrante:", w.registrar)
        print("  [-] Servidores de Nombres:", w.name_servers)
    except Exception as e:
        print(f"  [!] No se pudo obtener informacion WHOIS: {e}")

    print(f"\n[+] Realizando consultas DNS para: {domain}")
    
    record_types = ['A', 'AAAA', 'MX', 'NS', 'TXT', 'SOA']
    for record_type in record_types:
        try:
            answers = dns.resolver.resolve(domain, record_type)
            print(f"  [-] Registros {record_type}:")
            for rdata in answers:
                print(f"    {rdata.to_text()}")
        except (dns.resolver.NoAnswer, dns.resolver.NXDOMAIN):
            print(f"  [!] No se encontraron registros {record_type}.")
        except Exception as e:
            print(f"  [!] Error consultando registros {record_type}: {e}")

# ==============================================================================
# MODULO 2: ESCANER DE PUERTOS (BASE)
# (Este modulo se usara en tareas futuras)
# ==============================================================================
def scan_ports(target, port_spec):
    """
    Escanea los puertos especificados en el host objetivo.
    (Implementacion pendiente para futuras tareas)
    """
    print(f"[!] El modulo de escaneo de puertos para '{port_spec}' aun no esta implementado.")
    return []

\end{lstlisting}

\newpage

\section{Tarea para la Semana 1: Recopilación de Información Pasiva (DNS)}
\textbf{Tema General:} Obtener información de un dominio sin enviar paquetes directamente (más allá de las consultas DNS).

\subsection*{Objetivo Común}
El objetivo de esta semana es que cada estudiante tome la función base \texttt{get\_dns\_info} del archivo \texttt{modulos.py} y la descomponga en funciones más pequeñas y especializadas. Cada estudiante se centrará en un conjunto diferente de registros DNS.

\subsection*{Instrucciones para el Estudiante 1}
\begin{itemize}
    \item \textbf{Tarea:} Crear un script para obtener los registros \textbf{A} (direcciones IP) y \textbf{AAAA} (direcciones IPv6) de un dominio.
    \item \textbf{Pasos:}
    \begin{enumerate}
        \item En \texttt{modulos.py}, crea una nueva función llamada \texttt{get\_a\_records(domain)}.
        \item Adapta el código del bucle de la función \texttt{get\_dns\_info} para que tu nueva función consulte únicamente los tipos de registro 'A' y 'AAAA'.
        \item En \texttt{auditoria.py}, añade un nuevo argumento de línea de comandos, por ejemplo \texttt{--dns-a}, que al ser invocado, llame a tu nueva función \texttt{get\_a\_records}.
    \end{enumerate}
\end{itemize}

\subsection*{Instrucciones para el Estudiante 2}
\begin{itemize}
    \item \textbf{Tarea:} Crear un script para obtener los registros \textbf{MX} (servidores de correo) y \textbf{NS} (servidores de nombres) de un dominio.
    \item \textbf{Pasos:}
    \begin{enumerate}
        \item En \texttt{modulos.py}, crea una nueva función llamada \texttt{get\_mx\_ns\_records(domain)}.
        \item Adapta el código del bucle de la función \texttt{get\_dns\_info} para que tu nueva función consulte únicamente los tipos de registro 'MX' y 'NS'.
        \item En \texttt{auditoria.py}, añade un nuevo argumento de línea de comandos, por ejemplo \texttt{--dns-mxns}, que al ser invocado, llame a tu nueva función \texttt{get\_mx\_ns\_records}.
    \end{enumerate}
\end{itemize}

\subsection*{Instrucciones para el Estudiante 3}
\begin{itemize}
    \item \textbf{Tarea:} Crear un script para obtener los registros \textbf{TXT} (información descriptiva, como SPF) y \textbf{SOA} (inicio de autoridad) de un dominio.
    \item \textbf{Pasos:}
    \begin{enumerate}
        \item En \texttt{modulos.py}, crea una nueva función llamada \texttt{get\_txt\_soa\_records(domain)}.
        \item Adapta el código del bucle de la función \texttt{get\_dns\_info} para que tu nueva función consulte únicamente los tipos de registro 'TXT' y 'SOA'.
        \item En \texttt{auditoria.py}, añade un nuevo argumento de línea de comandos, por ejemplo \texttt{--dns-txtsoa}, que al ser invocado, llame a tu nueva función \texttt{get\_txt\_soa\_records}.
    \end{enumerate}
\end{itemize}

\subsection*{Nota Final}
Al finalizar la semana, el archivo \texttt{auditoria.py} debería poder invocar las funciones especializadas de cada estudiante. El objetivo es que aprendan a crear funciones modulares y a integrarlas en una aplicación principal mediante argumentos de línea de comandos.

\end{document}